% RESEARCH_DRAFT_SHA256: 47adce4512ea2bb8d4f4c81cbc1f9b2eb437cb238af3156bec9c7a2e95e3a93d
% Layout and source-reference wrappers expanded for the web renderer.
\documentclass[11pt]{article}
\usepackage[margin=25mm]{geometry}
\usepackage{amsthm}
\usepackage{fontspec,unicode-math}
\setmainfont{Latin Modern Roman}
\setmathfont{Latin Modern Math}
\usepackage{amsmath,amssymb,mathtools,xurl,hyperref}
\hypersetup{colorlinks=true,urlcolor=blue,linkcolor=blue}
\setcounter{secnumdepth}{0}
\setlength{\parindent}{0pt}
\setlength{\parskip}{5pt}
\setlength{\emergencystretch}{3em}
\newcommand{\sref}[2]{\hyperlink{source:#1:#2}{[S#2]}}
\newcommand{\cataloguescope}[1]{\paragraph{Source record.}#1}
\newcommand{\floor}[1]{\left\lfloor#1\right\rfloor}
\newtheorem{proposition}{Proposition}
\newtheorem{lemma}[proposition]{Lemma}
\begin{document}
\paragraph{Research manuscript.} OpenAI \texttt{gpt-6-astra}, reasoning \texttt{ultra}; 11 October 2026.
\paragraph{Preserved source body.} The complete catalogue mathematical text and sources follow with layout and source links expanded, including source qualifications and its historical preparation notes. New research begins at the section headed Research attempt.

% UnsolvedMath ID 20002103; source catalogue code AIM-INFRASTRUCTURE-0003
\section{A motivation-first deployment gate for educational technology}\label{20002103}
{\small UnsolvedMath ID 20002103 · Computer Science}
\subsection{Definitions and mathematical statement}

\paragraph{Shared notation.}Standard number systems are $\mathbb N,\mathbb Z,\mathbb Q,\mathbb R,\mathbb C$. Specialized notation and index ranges follow the cited source and the definitions below.


Student motivation concerns willingness to engage and persist in learning tasks. A study must distinguish a recorded behavior, such as participation, from the underlying motivational construct and define how either is measured.
\paragraph{Record scope and qualifications.}
The catalogue title is retained as metadata. Named certificates, contracts, bounds or frameworks in that title are not asserted as results of the original source extract.
This is a study topic without a population, intervention, outcome definition, or testable quantitative claim. The supplied catalogue marks this record open; its present research status was not independently checked.
The source text below is rendered literally where necessary; damaged formulas and composite extracts are not silently repaired into a new mathematical claim.
\paragraph{Statement recorded in the supplied source.}
3. Technology as it relates to how we motivate students \href{http://aimath.org/pastworkshops/oerassessmentproblems.pdf}{[S2]}
\paragraph{Source status.}
The supplied catalogue labels this source record open; that metadata was checked against the supplied snapshot on 2026-10-11. The label is preserved for provenance and is not an independent verification of present open status or mathematical correctness.
\subsection{Short English statement}
The source proposes studying the relation between technology and student motivation.
\subsection{Sources}
\begin{itemize}
\item[S1] Ulam AI and the UnsolvedMath Contributors, UnsolvedMath: A Curated Collection of Open Mathematics Problems, record 20002103 (AIM-INFRASTRUCTURE-0003). Catalogue curation and metadata: CC BY 4.0. The source extract is retained; definitions, English summaries and qualifications are editorial additions.
\url{https://huggingface.co/datasets/ulamai/UnsolvedMath}.
\item[S2] AIM Workshop Problem Lists, source linked by catalogue record AIM-INFRASTRUCTURE-0003.
\url{http://aimath.org/pastworkshops/oerassessmentproblems.pdf}.
\end{itemize}
\paragraph{Preparation and review.}
The original source and UnsolvedMath contributors retain attribution for the source extract. Codex (GPT-6; exact serving revision not exposed) assisted with definitions, English summaries, scope notes and formatting. No new human mathematical authorship or independently verified solution is asserted. Permanent source ID 20002103 and code AIM-INFRASTRUCTURE-0003 checked against the complete supplied registry and current repository. Distinct upstream records, including repeated questions and broad or damaged extracts, are retained for complete coverage of the source CS category. A repeated record is not a new independent conjecture. Complete mathematical-equivalence and current-literature review has not been performed. Preparation and final review follow contribution prompts C and D, with source gaps disclosed. Structural and rendering checks do not establish a new theorem or open status.
\label{end:20002103}

\section{Research attempt}
\subsection{Literature, source context and scope}
The primary AIM notes describe a brainstorming session. The original excerpt and all catalogue qualifications remain above. They do not specify a mathematical theorem or a complete empirical study. The formal model below is an explicitly limited research route introduced for this attempt, not a reconstruction of a missing original conjecture. The supplied catalogue's earlier research summary was inspected as context; its named framework is not adopted as an established solution. The elementary result below is not claimed novel, and no educational effectiveness data were collected.

\subsection{Attempt: incentives for recorded activity}
To examine the source's link between technology and motivation, I tested whether increasing the reward for a logged action necessarily selects the action that supports learning. This is a two-action incentive model, with learning outcome and utility deliberately separated; it is not a measured model of students.
\begin{proposition}
An activity reward can increase logged activity tenfold while changing the chosen learning outcome from one to zero.
\end{proposition}
\begin{proof}
Suppose exactly two actions are available. Action $L$ produces one logged unit and learning outcome one, with intrinsic value one and effort cost one. Action $C$ produces ten logged units and learning outcome zero, with no intrinsic value and cost $1/10$. Paying reward $r\ge0$ per logged unit gives utilities $U_L=r$ and $U_C=10r-1/10$. At $r=0$, $L$ uniquely beats $C$. For $r>1/90$, their difference $U_C-U_L=9r-1/10$ is positive, so $C$ is uniquely selected. The logged amount changes from one to ten while the stipulated learning outcome falls from one to zero.
\end{proof}
\subsection{Verification and remaining gaps}
At $r=1/10$, the utilities are $1/10$ and $9/10$. This exact comparison identifies the failed inference: rewarding a recorded proxy need not reward the desired action when the proxy can be produced by another route. The counterexample does not assert that real students behave this way, that extrinsic rewards always harm motivation, or that activity causes learning loss. Values, action sets and motivational constructs would need empirical specification before evaluating a deployment. The source is an open educational topic and no complete mathematical target is recovered or solved.

\section{References and source attribution}
American Institute of Mathematics, workshop discussion notes on open educational resources and assessment, \url{https://aimath.org/pastworkshops/oerassessmentproblems.pdf}. The 16-page primary PDF was opened on 11 October 2026 and its numbered brainstorming item was checked against the supplied record. The source supports the provenance and agenda, not the additional modeling assumptions or an empirical educational conclusion. Catalogue and contributor attribution are retained in the reproduced Sources section. The worked mathematics is proved here without reliance on an unverified external theorem.


\paragraph{Integration review.} Codex (GPT-6; exact serving revision not exposed) reviewed the written arguments and expanded layout and source-reference wrappers. Repository metadata and the production web renderer were checked using Python 3.13 and JavaScript with MathJax 3.2.2. This model-assisted review is not independent human proof verification; no PDF compilation or formal proof checking is asserted.

\section{Completion Estimate}
\noindent\textbf{COMPLETION: 0\%}
No substantiated new progress toward the entire original target is claimed. The worked results and failed route are delimited above; a known result, illustrative model, or auxiliary obstruction is not counted as original-target completion.

\section{Final status and attribution}
\textbf{Attempt status: unresolved.} This describes this research attempt, not a new ruling on the source catalogue's status or on external literature claims.
The mathematics and this self-audit were generated by OpenAI \texttt{gpt-6-astra}, reasoning \texttt{ultra}, on 11 October 2026. Codex (GPT-6; exact serving revision not exposed) orchestrated the batch and contributed formatting/repository review. No human mathematical authorship, independent human proof checking, or formal verification is claimed. The historical preparation paragraphs in the reproduced catalogue body concern that body only.

Contribution prompt A was applied to the complete definition and supplied source record, followed by the adversarial checks of prompt B. Local reproduction files are \texttt{batch008/checks.py}, \texttt{computations.json}, and the per-ID audit JSON. They were not public links at preparation time. Numerical checks are finite and do not prove unrestricted results. No LaTeX engine was found locally; compilation was not performed by this research agent.

\end{document}
